\section*{4. 树的节点着色（最大化黑色节点数）}

对一棵以r为根的树进行节点着色，每个节点可着黑色或白色。约束：相邻节点不能着相同黑色（但可同白）。求使树中黑色节点数最多的着色方案。

\begin{itemize}
    \item \textbf{状态定义}：$dp[u][0]$表示节点u着白色时子树最大黑色节点数，$dp[u][1]$表示节点u着黑色时子树最大黑色节点数
    \item \textbf{转移方程}：$dp[u][1] = 1 + \sum_{v \in children(u)} dp[v][0]$，$dp[u][0] = \sum_{v \in children(u)} \max(dp[v][0], dp[v][1])$
    \item \textbf{遍历顺序}：后序遍历（先子节点后父节点）
    \item \textbf{最终结果}：$\max(dp[r][0], dp[r][1])$
\end{itemize}

\begin{lstlisting}[language=C++, style=cppstyle]
struct TreeNode {
    vector<TreeNode*> children;
};

pair<int, int> treeColoring(TreeNode* u) {
    // 返回 {着白时的最大黑数, 着黑时的最大黑数}
    int white = 0, black = 1;  // 初始化：无子节点时

    for (TreeNode* v : u->children) {
        auto [w, b] = treeColoring(v);

        // u着黑：子节点必须着白
        black += w;

        // u着白：子节点可选最大
        white += max(w, b);
    }

    return {white, black};
}

int maxBlackNodes(TreeNode* root) {
    auto [white, black] = treeColoring(root);
    return max(white, black);
}
\end{lstlisting}

\begin{enumerate}
    \item 从根节点开始后序遍历（先处理完所有孩子，再处理自己）
    \item 对每个节点算两个值：染白时子树最多几个黑节点、染黑时子树最多几个黑节点
    \item 若染黑，所有孩子只能染白；若染白，每个孩子自由选择
    \item 根节点取两者较大值
\end{enumerate}
